% Options for packages loaded elsewhere
\PassOptionsToPackage{unicode}{hyperref}
\PassOptionsToPackage{hyphens}{url}
\documentclass[
]{article}
\usepackage{xcolor}
\usepackage{amsmath,amssymb}
\setcounter{secnumdepth}{-\maxdimen} % remove section numbering
\usepackage{iftex}
\ifPDFTeX
  \usepackage[T1]{fontenc}
  \usepackage[utf8]{inputenc}
  \usepackage{textcomp} % provide euro and other symbols
\else % if luatex or xetex
  \usepackage{unicode-math} % this also loads fontspec
  \defaultfontfeatures{Scale=MatchLowercase}
  \defaultfontfeatures[\rmfamily]{Ligatures=TeX,Scale=1}
\fi
\usepackage{lmodern}
\ifPDFTeX\else
  % xetex/luatex font selection
\fi
% Use upquote if available, for straight quotes in verbatim environments
\IfFileExists{upquote.sty}{\usepackage{upquote}}{}
\IfFileExists{microtype.sty}{% use microtype if available
  \usepackage[]{microtype}
  \UseMicrotypeSet[protrusion]{basicmath} % disable protrusion for tt fonts
}{}
\makeatletter
\@ifundefined{KOMAClassName}{% if non-KOMA class
  \IfFileExists{parskip.sty}{%
    \usepackage{parskip}
  }{% else
    \setlength{\parindent}{0pt}
    \setlength{\parskip}{6pt plus 2pt minus 1pt}}
}{% if KOMA class
  \KOMAoptions{parskip=half}}
\makeatother
\usepackage{color}
\usepackage{fancyvrb}
\newcommand{\VerbBar}{|}
\newcommand{\VERB}{\Verb[commandchars=\\\{\}]}
\DefineVerbatimEnvironment{Highlighting}{Verbatim}{commandchars=\\\{\}}
% Add ',fontsize=\small' for more characters per line
\newenvironment{Shaded}{}{}
\newcommand{\AlertTok}[1]{\textcolor[rgb]{1.00,0.00,0.00}{\textbf{#1}}}
\newcommand{\AnnotationTok}[1]{\textcolor[rgb]{0.38,0.63,0.69}{\textbf{\textit{#1}}}}
\newcommand{\AttributeTok}[1]{\textcolor[rgb]{0.49,0.56,0.16}{#1}}
\newcommand{\BaseNTok}[1]{\textcolor[rgb]{0.25,0.63,0.44}{#1}}
\newcommand{\BuiltInTok}[1]{\textcolor[rgb]{0.00,0.50,0.00}{#1}}
\newcommand{\CharTok}[1]{\textcolor[rgb]{0.25,0.44,0.63}{#1}}
\newcommand{\CommentTok}[1]{\textcolor[rgb]{0.38,0.63,0.69}{\textit{#1}}}
\newcommand{\CommentVarTok}[1]{\textcolor[rgb]{0.38,0.63,0.69}{\textbf{\textit{#1}}}}
\newcommand{\ConstantTok}[1]{\textcolor[rgb]{0.53,0.00,0.00}{#1}}
\newcommand{\ControlFlowTok}[1]{\textcolor[rgb]{0.00,0.44,0.13}{\textbf{#1}}}
\newcommand{\DataTypeTok}[1]{\textcolor[rgb]{0.56,0.13,0.00}{#1}}
\newcommand{\DecValTok}[1]{\textcolor[rgb]{0.25,0.63,0.44}{#1}}
\newcommand{\DocumentationTok}[1]{\textcolor[rgb]{0.73,0.13,0.13}{\textit{#1}}}
\newcommand{\ErrorTok}[1]{\textcolor[rgb]{1.00,0.00,0.00}{\textbf{#1}}}
\newcommand{\ExtensionTok}[1]{#1}
\newcommand{\FloatTok}[1]{\textcolor[rgb]{0.25,0.63,0.44}{#1}}
\newcommand{\FunctionTok}[1]{\textcolor[rgb]{0.02,0.16,0.49}{#1}}
\newcommand{\ImportTok}[1]{\textcolor[rgb]{0.00,0.50,0.00}{\textbf{#1}}}
\newcommand{\InformationTok}[1]{\textcolor[rgb]{0.38,0.63,0.69}{\textbf{\textit{#1}}}}
\newcommand{\KeywordTok}[1]{\textcolor[rgb]{0.00,0.44,0.13}{\textbf{#1}}}
\newcommand{\NormalTok}[1]{#1}
\newcommand{\OperatorTok}[1]{\textcolor[rgb]{0.40,0.40,0.40}{#1}}
\newcommand{\OtherTok}[1]{\textcolor[rgb]{0.00,0.44,0.13}{#1}}
\newcommand{\PreprocessorTok}[1]{\textcolor[rgb]{0.74,0.48,0.00}{#1}}
\newcommand{\RegionMarkerTok}[1]{#1}
\newcommand{\SpecialCharTok}[1]{\textcolor[rgb]{0.25,0.44,0.63}{#1}}
\newcommand{\SpecialStringTok}[1]{\textcolor[rgb]{0.73,0.40,0.53}{#1}}
\newcommand{\StringTok}[1]{\textcolor[rgb]{0.25,0.44,0.63}{#1}}
\newcommand{\VariableTok}[1]{\textcolor[rgb]{0.10,0.09,0.49}{#1}}
\newcommand{\VerbatimStringTok}[1]{\textcolor[rgb]{0.25,0.44,0.63}{#1}}
\newcommand{\WarningTok}[1]{\textcolor[rgb]{0.38,0.63,0.69}{\textbf{\textit{#1}}}}
\usepackage{longtable,booktabs,array}
\newcounter{none} % for unnumbered tables
\usepackage{calc} % for calculating minipage widths
% Correct order of tables after \paragraph or \subparagraph
\usepackage{etoolbox}
\makeatletter
\patchcmd\longtable{\par}{\if@noskipsec\mbox{}\fi\par}{}{}
\makeatother
% Allow footnotes in longtable head/foot
\IfFileExists{footnotehyper.sty}{\usepackage{footnotehyper}}{\usepackage{footnote}}
\makesavenoteenv{longtable}
\setlength{\emergencystretch}{3em} % prevent overfull lines
\providecommand{\tightlist}{%
  \setlength{\itemsep}{0pt}\setlength{\parskip}{0pt}}
\usepackage{bookmark}
\IfFileExists{xurl.sty}{\usepackage{xurl}}{} % add URL line breaks if available
\urlstyle{same}
\hypersetup{
  hidelinks,
  pdfcreator={LaTeX via pandoc}}

\author{}
\date{}

\begin{document}

\section{Trinity GF16: A phi-Anchored 16-bit Float with FPGA
Synthesised with an Open-Source FPGA
Toolchain}\label{trinity-gf16-a-phi-anchored-16-bit-float-synthesised-with-an-open-source-fpga-toolchain}

\subsection{Abstract}\label{abstract}

We introduce Golden Float 16 (GF16), a 16-bit floating-point format with
a phi-anchored exponent bias of 31. GF16 uses a 1/6/9 bit layout
(sign/exponent/mantissa) and is synthesised (no operating frequency is claimed; see Timing)
on a Xilinx Artix-7 XC7A100T FPGA using the open-source openXC7
toolchain (Yosys + nextpnr). We present a complete dot-product (N=4) and
4x4 matrix multiplication accelerator verified in FPGA synthesis and
RTL simulation, with 35/35 tests passing and 0 timing violations at
100 MHz. The design has been submitted for ASIC fabrication on Sky130
via the TinyTapeout TTSKY26b TT4913 Gamma shuttle (submission closed
May 2026); silicon has not yet been returned (expected late 2026), so
no on-chip measurement is claimed.

\subsection{1. Introduction}\label{introduction}

The golden ratio phi (phi = (1+sqrt(5))/2 = 1.618\ldots) appears
throughout nature, mathematics, and information theory. We observe that
the IEEE 754 half-precision format (float16) uses a 5-bit exponent with
bias 15, while bfloat16 uses 8-bit exponent with bias 127. Neither has
any connection to fundamental mathematical constants.

Trinity GF16 anchors its exponent bias at 31, derived from the
relationship:

\begin{verbatim}
phi^2 + phi^-2 = 3
phi^2 = phi + 1
\end{verbatim}

The bias value 31 encodes 1.0 at the golden ratio's ``natural center''
of the exponent range, creating a format where common physics and ML
values cluster around the representational sweet spot.

\subsection{2. GF16 Format
Specification}\label{gf16-format-specification}

\subsubsection{2.1 Bit Layout}\label{bit-layout}

{\def\LTcaptype{none} % do not increment counter
\begin{longtable}[]{@{}lll@{}}
\toprule\noalign{}
Bit(s) & Field & Width \\
\midrule\noalign{}
\endhead
\bottomrule\noalign{}
\endlastfoot
15 & Sign & 1 \\
14:9 & Exponent & 6 \\
8:0 & Mantissa & 9 \\
\end{longtable}
}

\begin{itemize}
\tightlist
\item
  \textbf{Total:} 16 bits
\item
  \textbf{Exponent bias:} 31
\item
  \textbf{Implicit leading 1:} Normal numbers have implicit 1.mantissa
\end{itemize}

\subsubsection{2.2 Special Values}\label{special-values}

{\def\LTcaptype{none} % do not increment counter
\begin{longtable}[]{@{}ll@{}}
\toprule\noalign{}
Value & Encoding \\
\midrule\noalign{}
\endhead
\bottomrule\noalign{}
\endlastfoot
+0 & 0x0000 \\
-0 & 0x8000 \\
+Infinity & 0x7E00 \\
-Infinity & 0xFE00 \\
NaN & 0xFE01 \\
\end{longtable}
}

\subsubsection{2.3 Encoding of Key
Constants}\label{encoding-of-key-constants}

{\def\LTcaptype{none} % do not increment counter
\begin{longtable}[]{@{}llll@{}}
\toprule\noalign{}
Value & GF16 Hex & Decoded & Relative Error \\
\midrule\noalign{}
\endhead
\bottomrule\noalign{}
\endlastfoot
1.0 & 0x3E00 & 1.0 & 0 \\
phi & 0x3F3C & 1.6171875 & 0.0005 \\
pi & 0x4049 & 3.140625 & 0.0003 \\
e & 0x4058 & 2.71875 & 0.0002 \\
sqrt(2) & 0x3F5C & 1.4140625 & 0.0001 \\
\end{longtable}
}

\subsubsection{2.4 Dynamic Range}\label{dynamic-range}

\begin{itemize}
\tightlist
\item
  \textbf{Max normal:} (1 + 511/512) x 2\^{}(62-31) =
  \textasciitilde4.29 x 10\^{}9
\item
  \textbf{Min normal:} 1.0 x 2\^{}(1-31) = \textasciitilde9.31 x
  10\^{}-10
\item
  \textbf{Machine epsilon:} 2\^{}-9 = 0.001953125
\end{itemize}

\subsubsection{2.5 Comparison with Existing
Formats}\label{comparison-with-existing-formats}

{\def\LTcaptype{none} % do not increment counter
\begin{longtable}[]{@{}llllll@{}}
\toprule\noalign{}
Format & Bits & Exp & Mant & Bias & Max Value \\
\midrule\noalign{}
\endhead
\bottomrule\noalign{}
\endlastfoot
float16 & 16 & 5 & 10 & 15 & 65504 \\
\textbf{GF16} & \textbf{16} & \textbf{6} & \textbf{9} & \textbf{31} &
\textbf{4.29 x 10\^{}9} \\
bfloat16 & 16 & 8 & 7 & 127 & 3.39 x 10\^{}38 \\
\end{longtable}
}

GF16 provides 65x wider dynamic range than float16 while maintaining
better precision than bfloat16 (9 vs 7 mantissa bits). The 6-bit
exponent with bias=31 positions 1.0 at the center of a practical
ML/inference range.

\subsection{3. Hardware Architecture}\label{hardware-architecture}

\subsubsection{3.1 GF16 Multiplier}\label{gf16-multiplier}

Combinational multiplier implementing: - 10-bit x 10-bit mantissa
multiplication (mapped to DSP48E1 on Xilinx) - Exponent addition with
bias subtraction - Round-to-nearest-even with guard/round/sticky bits -
Special value handling (NaN, Inf, zero)

\subsubsection{3.2 GF16 Adder}\label{gf16-adder}

Combinational adder implementing: - Exponent alignment via priority
encoder (no for-loops) - Sign-magnitude addition with cancellation
support - Normalization with round-to-nearest-even - Special value
handling

Key design constraint: all internal \texttt{reg} variables initialized
with defaults at the top of \texttt{always\ @(*)} blocks to prevent
latch inference in Yosys synthesis.

\subsubsection{3.3 Dot Product (N=4)}\label{dot-product-n4}

Tree reduction: 4 multipliers + 3 adders

\begin{verbatim}
a0*b0 ---\
          +-- s01 ---\
a1*b1 ---/           \
                      +-- result
a2*b2 ---\           /
          +-- s23 --/
a3*b3 ---/
\end{verbatim}

\subsubsection{3.4 Matrix Multiplication
(4x4)}\label{matrix-multiplication-4x4}

16 parallel dot-product units. Each unit computes one element of the
output matrix C = A x B.

\begin{verbatim}
C[i][j] = dot4(A[i][0:3], B[0:3][j])
\end{verbatim}

\subsection{4. FPGA Implementation
Results}\label{fpga-implementation-results}

\subsubsection{4.1 Platform}\label{platform}

\begin{itemize}
\tightlist
\item
  \textbf{FPGA:} QMTECH XC7A100T-1FGG676C (Artix-7, 126,800 LUTs, 240
  DSP48E1)
\item
  \textbf{Toolchain:} openXC7 (Yosys 0.62 + nextpnr-xilinx)
\item
  \textbf{Clock:} 20-stage ring oscillator (onboard M21 crystal
  non-functional)
\item
  \textbf{Programming:} XVC via ESP32 (192.168.1.30:2542)
\end{itemize}

\subsubsection{4.2 Resource Utilization}\label{resource-utilization}

{\def\LTcaptype{none} % do not increment counter
\begin{longtable}[]{@{}lllll@{}}
\toprule\noalign{}
Design & LUTs & DSP48E1 & Tests \\
\midrule\noalign{}
\endhead
\bottomrule\noalign{}
\endlastfoot
GF16 mul & \textasciitilde650 & 1 & 13/13 \\
GF16 dot4 & 2,605 & 4 & 6/6 \\
GF16 matmul 4x4 & 40,350 & 64 & 35/35 \\
\end{longtable}
}

\subsubsection{4.3 Timing}\label{timing}

Resource figures are \texttt{yosys stat} on the DUT. They are \textbf{not} from the placed wrapper that produced the bitstream: that wrapper feeds the DUT literal constants, so the arithmetic is constant-folded out and its synthesised design module holds 55 cells and no DSP48E1. A \texttt{Max Freq} column stood here and is withdrawn --- all three of its entries came off the same ring-oscillator net:

\begin{verbatim}
[withdrawn: this clock is a ring-oscillator probe, not a GF16 path]
\end{verbatim}

\subsubsection{4.4 Latch Elimination}\label{latch-elimination}

A critical design challenge was eliminating all LDCE latch inferences
from Yosys. The root cause was incompletely-assigned \texttt{reg}
variables in \texttt{always\ @(*)} combinational blocks. Solution:
initialize all \texttt{reg} variables with defaults at the block entry,
and ensure every control path assigns every variable.

Before fix: 1 latch (gf16\_add.result) After fix: 0 latches, 0 errors

\subsubsection{4.5 Hardware Verification}\label{hardware-verification}

All designs verified on FPGA via XVC programming: - dot4({[}1,2,3,4{]},
{[}1,2,3,4{]}) = 30.0 (0x47C0) --- LEDs confirm - matmul4x4 identity
test --- ISC\_DONE=1, DONE=1

\subsubsection{4.6 Throughput}\label{throughput}

{\def\LTcaptype{none} % do not increment counter
\begin{longtable}[]{@{}ll@{}}
\toprule\noalign{}
Metric & Value \\
\midrule\noalign{}
\endhead
\bottomrule\noalign{}
\endlastfoot
GF16 ops per matmul4x4 & 128 (64 mul + 64 add) \\
Throughput at a clock of $f$ & $128f$ ops/sec, fully parallel, 1 cycle \\
\end{longtable}
}

\subsection{5. ASIC Path (TinyTapeout
TTSKY26b TT4913 Gamma)}\label{asic-path-tinytapeout-ttsky26b}

\subsubsection{5.1 Submission}\label{submission}

The GF16 dot4 design was wrapped for the TinyTapeout Sky130 shuttle:

\begin{itemize}
\tightlist
\item
  \textbf{Repo:} github.com/gHashTag/tt-trinity-gf16
\item
  \textbf{Module:} tt\_um\_ghtag\_trinity\_gf16
\item
  \textbf{Tile:} 1x1 (167 x 108 um)
\item
  \textbf{PDK:} Sky130A (OpenLane/LibreLane 3.0.0)
\end{itemize}

\subsubsection{5.2 CI Results}\label{ci-results}

{\def\LTcaptype{none} % do not increment counter
\begin{longtable}[]{@{}ll@{}}
\toprule\noalign{}
Check & Status \\
\midrule\noalign{}
\endhead
\bottomrule\noalign{}
\endlastfoot
GDS Build & PASSED \\
Gate-Level Test & PASSED \\
Precheck & PASSED \\
DRC & Clean \\
LVS & Clean \\
\end{longtable}
}

\subsubsection{5.3 Expected Timeline}\label{expected-timeline}

\begin{itemize}
\tightlist
\item
  Shuttle close: May 11, 2026
\item
  Chip delivery: \textasciitilde December 2026
\end{itemize}

\subsection{6. Python Reference
Implementation}\label{python-reference-implementation}

A complete Python reference (encode/decode/mul/add/dot4) is provided in
\texttt{conformance/gf16\_ref.py}:

\begin{itemize}
\tightlist
\item
  32/32 FPGA consistency tests pass
\item
  19/19 roundtrip tests pass
\item
  Encoding matches FPGA testbench exactly
\end{itemize}

\subsection{7. Conclusion}\label{conclusion}

We have demonstrated a complete implementation of the Trinity GF16
floating-point format, from specification through FPGA verification to
ASIC submission. The phi-anchored bias=31 provides a natural centering
for ML and scientific computation values, while the 6/9
exponent/mantissa split offers 65x wider dynamic range than float16 with
better precision than bfloat16.

What this work establishes, and no more: the design synthesises for the
Artix-7 XC7A100T through an entirely open-source flow, its testbench passes
35/35, and Yosys infers no latches. The resource figures are synthesis
statistics on the DUT. \textbf{No operating frequency, and therefore no
throughput, is claimed} --- the figure previously reported was a
ring-oscillator probe and is withdrawn (see Timing). No ASIC silicon has been
measured.

\subsection{References}\label{references}

\begin{enumerate}
\def\labelenumi{\arabic{enumi}.}
\tightlist
\item
  IEEE 754-2019, IEEE Standard for Floating-Point Arithmetic
\item
  TinyTapeout, https://tinytapeout.com
\item
  openXC7 toolchain, https://github.com/openXC7
\item
  Yosys Open SYnthesis Suite, https://github.com/YosysHQ/yosys
\item
  nextpnr-xilinx, https://github.com/openXC7/nextpnr-xilinx
\end{enumerate}

\subsection{Appendix A: GF16 Encoding
Examples}\label{appendix-a-gf16-encoding-examples}

\begin{verbatim}
Value    | Hex    | Sign | Exp | Mant | Decoded
---------|--------|------|-----|------|--------
0.0      | 0x0000 | 0    | 0   | 0    | 0.0
1.0      | 0x3E00 | 0    | 31  | 0    | 1.0
-1.0     | 0xBE00 | 1    | 31  | 0    | -1.0
2.0      | 0x4000 | 0    | 32  | 0    | 2.0
0.5      | 0x3C00 | 0    | 30  | 0    | 0.5
3.0      | 0x4100 | 0    | 32  | 256  | 3.0
1.5      | 0x3F00 | 0    | 31  | 256  | 1.5
100.0    | 0x4B20 | 0    | 37  | 288  | 100.0
phi      | 0x3F3C | 0    | 31  | 60   | 1.6171875
+Inf     | 0x7E00 | 0    | 63  | 0    | Infinity
NaN      | 0xFE01 | 1    | 63  | 1    | NaN
\end{verbatim}

\subsection{Appendix B: Reproduce
Instructions}\label{appendix-b-reproduce-instructions}

\begin{Shaded}
\begin{Highlighting}[]
\CommentTok{\# FPGA synthesis (openXC7)}
\ExtensionTok{docker}\NormalTok{ run }\AttributeTok{{-}{-}rm} \AttributeTok{{-}v} \StringTok{"}\VariableTok{$(}\BuiltInTok{pwd}\VariableTok{)}\StringTok{"}\NormalTok{:/workspace regymm/openxc7 bash }\AttributeTok{{-}c} \StringTok{\textquotesingle{}}
\StringTok{  yosys {-}p "read\_verilog fpga/vivado/gf16\_*.v; synth\_xilinx {-}top gf16\_matmul4x4 {-}family xc7; write\_json build/synth.json"}
\StringTok{  nextpnr{-}xilinx {-}{-}chipdb xc7a100t.bin {-}{-}json build/synth.json {-}{-}freq 100}
\StringTok{\textquotesingle{}}

\CommentTok{\# RTL simulation}
\ExtensionTok{iverilog} \AttributeTok{{-}g2012} \AttributeTok{{-}o}\NormalTok{ /tmp/test fpga/vivado/gf16\_}\PreprocessorTok{*}\NormalTok{.v }\KeywordTok{\&\&} \ExtensionTok{vvp}\NormalTok{ /tmp/test}

\CommentTok{\# Python reference}
\ExtensionTok{python3}\NormalTok{ conformance/gf16\_ref.py}
\end{Highlighting}
\end{Shaded}


\end{document}
